\chapter{多模态模型}
\label{chap:vision-multimodal-models}

同一张街景照片，可以用来寻找相似图片、回答“右侧停着什么车”，也可以作为生成另一种天气下街景的条件。若再给一段录音，系统还能比较画面与声音是否相关，或把回答说出来。这些功能表面上都接受多种输入，内部需要的计算却不同。一个能比较图文相似度的模型，未必有生成句子的输出层；一个会说话的视觉助手，也未必能生成视频。

前六章围绕视觉成果介绍了具体任务。本章转向这些任务共用的模型接口，追踪观测怎样被编码、不同来源的信息在哪里相遇、输出又怎样恢复成分数、语言或媒体。编码和连接决定模型能利用哪些证据，训练目标决定哪些对应关系得到学习，输出模块决定结果能够以什么形式出现。理解和生成可以共享其中一部分，也可以由几个专用模型协作完成。

\section{多模态模型的组成与接口}
\label{sec:vision-models-interfaces}

连接图像和文字之前，需要区分原始观测、网络内部表示和最终结果。它们的数值形式可能相似，却承担不同职责。例如，一张照片可以形成许多空间位置上的特征，也可以被汇总为一个整图向量；后者便于匹配，却未必足以恢复前者包含的小字和边界。

\subsection{各模态的编码与组织}
\label{sec:vision-models-encoding}

沿用前文记号，图像为$\bm{X}\in\mathbb{R}^{h\times w\times c}$，视频由带时间戳的帧$\bm{X}_t$组成。图像编码器把像素转换为$\bm{F}\in\mathbb{R}^{n_{\mathrm v}\times d_{\mathrm v}}$，每一行表示一个位置或一个汇总单元。对图像块编码，$n_{\mathrm v}$通常随空间分辨率增加。位置编码使模型能够区分“左侧的车”和“右侧的车”，但位置编码不能补回裁剪时删除的对象。

文字先被切成离散词元，再经嵌入表转为向量。音频则可以由波形转换为时频表示，再划分成局部单元编码。时频表示的两维分别描述时间和频率，和图像的横纵坐标不同。视频还需保留帧的实际时间、帧内空间位置和所属片段；单纯把第1帧、第2帧写入序列，无法区分两帧相隔0.1秒还是10秒。

实际系统常把连续的视觉特征也称为“视觉词元”。这里的词元表示可进入序列模型的一个计算单元，不一定经过离散词表量化。它与文字词元的整数编号、媒体生成器使用的离散编码、原始像素值，都不能直接互换。比较词元数时，必须说明每个单元对应多少空间或时间范围。

设街景原图宽1200像素、高800像素，裁出左上角为$(400,200)$、宽高为$400\times400$的区域，再调整到$200\times200$作为局部输入。以连续边界坐标计，局部结果$(x',y')=(50,80)$对应原图
\[
  (x,y)=(400+2x',\,200+2y')=(500,360)\text{。}
\]
因此，输入记录至少应保留原图编号、裁剪边界和缩放关系。若视频只选取0、2、4秒三帧，还应保留真实时间戳，而不是让下游把三帧当成连续拍摄的相邻帧。编码前的选择已经限定了模型可能看到的证据。

\subsection{跨模态信息的连接}
\label{sec:vision-models-connection}

\term{视觉语言模型}（vision-language model, VLM）把视觉输入与语言联系起来，可以输出可比较表示，也可以生成语言。本书采用这一较宽的用法。连接音频、深度或惯性信号的模型还需相应编码和训练依据，不能仅凭“多模态”名称确定支持范围。

一种安排是两侧独立编码，再把表示映射到共同空间计算相似度。图像向量不依赖当前文字查询，适合预先建立图库。另一种安排让文字计算读取视觉特征，例如以语言隐状态为查询、视觉特征为键和值执行交叉注意力。此时视觉底层特征可以缓存，但读取结果随问题和已生成文字改变。还有一种安排把多源表示放入共享主干，通过位置、模态标记和注意力范围组织交互。

这三种安排描述计算发生在哪里，不是互斥的完整模型类别。一个系统可以先用独立表示召回候选，再让联合模型逐个检查。表示对齐指训练使不同来源的内容在比较或交互中形成有用联系，参数共享指不同输入实际使用相同的权重，连接模块则负责把形状或信息路径接起来。两套编码器即使输出同样维数，也没有因此完成语义对齐。

以图库查询为例，独立编码可以保存$M$幅图的向量，新请求只需编码一次文字并比较候选。若每幅图都要与当前问题共同通过一个大型主干，通常就要执行$M$次相关联合计算。前者能否找到关系复杂的目标、后者能否在预算内处理全部候选，都需要按任务确认；模块安排只说明这些成本和能力从何而来。

\subsection{输出表示与解码}
\label{sec:vision-models-output}

输出分数、输出文字和输出媒体需要不同的预测对象。匹配接口可以计算一个向量内积；语言接口在每一步给出词表上的分布；图像或音频接口则预测像素、波形或能够被解码的中间表示。中间表示必须与解码器的训练约定相容，任意同维向量都不能送入同一个解码器。

表\ref{tab:vision-model-interfaces}固定同一街景输入，比较不同请求所需的结果。若给CLIP附加图像生成器，CLIP可以提供条件，但生成能力属于完整连接后的系统。反过来，视觉问答虽然输出文字，评价的仍是回答是否依据画面作出了正确判断。

\begin{table}[htbp]
  \centering
  \small
  \begin{tabularx}{\linewidth}{@{}p{22mm}XX@{}}
    \toprule
    请求 & 模型直接预测的对象 & 恢复最终结果还需什么\\
    \midrule
    选择车辆类别 & 候选描述的匹配分数 & 给定候选集与选择规则\\
    回答停车位置 & 下一文字词元的分布 & 逐词解码与来源核验\\
    生成雨天变体 & 图像潜表示的更新量 & 生成迭代与图像解码器\\
    说出判断 & 语音编码序列 & 声学解码与波形合成\\
    \bottomrule
  \end{tabularx}
  \caption{输入相同而输出接口不同的教学比较。每行描述一种完整路径，并不意味着同一个检查点同时具有四种能力。}
  \label{tab:vision-model-interfaces}
\end{table}
\FloatBarrier

\section{多模态理解模型}
\label{sec:vision-models-understanding}

如果任务是从给定候选中选择，模型可以用可比较表示形成判断。如果任务要求自由描述或回答，则还需利用语言接口组织结果。本节先展开匹配如何得到训练，再说明这种接口的用途与边界；语言响应的具体连接方式在\ref{sec:vision-models-language-generation}节展开。

\subsection{CLIP的图文双编码与相似度判定}
\label{sec:vision-models-contrastive}

为了用文字指定视觉概念，需要让“红色汽车”这样的描述与包含相应内容的图片能够比较。CLIP以真实图文配对为监督，分别编码两种输入，再把正确配对与同一批次中的其他组合区分开。其图像编码器可以采用经修改的ResNet或ViT，文字编码器采用Transformer，双编码结构不要求两边网络相同\scite{radford2021}

设一批有$n$对图文$(\bm{X}_i,q_i)$。图像和文字编码器分别输出$f_{\mathrm v}(\bm{X}_i)$与$f_{\mathrm t}(q_i)$，线性投影矩阵$\bm{W}_{\mathrm v}$、$\bm{W}_{\mathrm t}$将它们映射到$\mathbb{R}^d$。对于非零投影结果，归一化表示为
\begin{equation}
  \bm{v}_i=\frac{\bm{W}_{\mathrm v}f_{\mathrm v}(\bm{X}_i)}
                    {\lVert\bm{W}_{\mathrm v}f_{\mathrm v}(\bm{X}_i)\rVert_2},
  \qquad
  \bm{u}_i=\frac{\bm{W}_{\mathrm t}f_{\mathrm t}(q_i)}
                    {\lVert\bm{W}_{\mathrm t}f_{\mathrm t}(q_i)\rVert_2}\text{。}
  \label{eq:vision-clip-embedding}
\end{equation}
内积$s_{i,j}=\bm{v}_i^\top\bm{u}_j$就是余弦相似度。归一化使向量长度不能单独抬高分数，模型必须改变方向来改善配对。原CLIP共同学习编码器、线性投影以及控制分数尺度的温度参数。

温度$\tau>0$把相似度转换为$\ell_{i,j}=s_{i,j}/\tau$。固定一幅图，在本批次$n$段文字中识别其配对文本，可以使用行方向交叉熵；固定一段文字，反过来在图像中识别配对，则使用列方向交叉熵：
\begin{equation}
  \begin{aligned}
    L_{\mathrm{v\to t}}
      &=-\frac1n\sum_{i=1}^n
        \log\frac{\exp(\ell_{i,i})}{\sum_{j=1}^n\exp(\ell_{i,j})},\\
    L_{\mathrm{t\to v}}
      &=-\frac1n\sum_{i=1}^n
        \log\frac{\exp(\ell_{i,i})}{\sum_{j=1}^n\exp(\ell_{j,i})},\\
    L_{\mathrm{CLIP}}&=\tfrac12
        \bigl(L_{\mathrm{v\to t}}+L_{\mathrm{t\to v}}\bigr)\text{。}
  \end{aligned}
  \label{eq:vision-clip-loss}
\end{equation}
对角线表示数据记录中的正确配对，其他单元在该目标中充当负例。两项损失约束相同的相似度矩阵，但分母分别沿行和列求和。若批次中另有一幅也符合该描述的图，单一配对标签会把它当负例，这就是配对监督和真实语义关系之间的差别。

取一个可复算的二维教学批次：两幅图的表示分别为$(1,0)^\top$、$(0,1)^\top$，两段配对文字的表示分别为$(0.8,0.6)^\top$、$(0.6,0.8)^\top$。四个向量均为单位向量。设两组标签为“红色汽车”和“红色自行车”，并取$\tau=0.2$，就得到
\[
  \bm{S}=\begin{bmatrix}0.8&0.6\\0.6&0.8\end{bmatrix},
  \qquad
  \bm{\ell}=\begin{bmatrix}4&3\\3&4\end{bmatrix}\text{。}
\]
每一行的正确配对概率为$e^4/(e^4+e^3)\approx0.7311$，交叉熵为$\log(1+e^{-1})\approx0.3133$；列方向在这个对称例子中相同。行损失对第一行两个logit的梯度，暂不计批次平均系数，为$(-0.2689,0.2689)$。梯度下降会提高正确配对相对错误配对的分数。图\ref{fig:vision-clip-matrix}标出了两个归一化方向。

\begin{figure}[htbp]
  \centering
  % 教学数据由四个单位向量的内积构造，温度为0.2。
\begin{tikzpicture}[x=1cm,y=1cm,font=\figurefont\small,
  arr/.style={-{Stealth[length=2mm]},BookInk,line width=0.8pt}]
  \node[anchor=west] at (0,4.5) {缩放后分数$\ell_{i,j}=s_{i,j}/\tau$};
  \node[align=center] at (3.8,3.55) {$q_1$\\红色汽车};
  \node[align=center] at (6.0,3.55) {$q_2$\\红色自行车};
  \node[anchor=east] at (2.6,2.5) {图像1};
  \node[anchor=east] at (2.6,1.2) {图像2};
  \foreach \x/\y/\val in {3.8/2.5/4,6.0/2.5/3,3.8/1.2/3,6.0/1.2/4}{
    \draw[BookRule,line width=0.5pt,fill=BookPaper]
      (\x-1,\y-0.55) rectangle (\x+1,\y+0.55);
    \node at (\x,\y) {$\val$};
  }
  \draw[BookTeal,line width=1.2pt] (2.8,1.95) rectangle (4.8,3.05);
  \draw[BookTeal,line width=1.2pt] (5.0,0.65) rectangle (7.0,1.75);
  \draw[arr] (7.2,2.5) -- (8.15,2.5);
  \node[anchor=west,align=left] at (8.2,2.5) {行softmax\\$0.7311$\\$0.2689$};
  \draw[arr] (3.8,0.4) -- (3.8,-0.2);
  \node[align=center,anchor=north] at (3.8,-0.3)
    {列softmax\\$(0.7311,\;0.2689)^\top$};
  \node[anchor=west] at (0,-1.35) {对角线：数据配对\qquad 非对角线：批次负例};
\end{tikzpicture}

  \caption{CLIP双向目标的两对图文教学示例。行固定图像、列固定文字，粗框标出训练配对。行、列softmax分别形成两个分类问题，图中的0.7311是批次内选择概率，不是“图中确有该物体”的校准概率。}
  \label{fig:vision-clip-matrix}
\end{figure}

使用时，可以把类别$c$写成描述模板$q_c$，缓存其表示$\bm{u}_c$，再对新图像选取$\argmax_c\bm{v}^\top\bm{u}_c$。写“这是一张红色汽车的照片”和只写“汽车”，提供给文字编码器的上下文不同，得到的方向也可能不同。若任务只有“汽车”“自行车”两个候选，模型一定会从中选一个，除非另加拒绝规则；加入“公交车”后，归一化分数的分母又会改变。

这种零样本分类不使用当前类别任务的标注样本去拟合分类器，但仍依赖大规模图文预训练。比较结果时应记录候选集合、模板、图像处理和模型版本。整图相似度没有直接给出车辆的框、实例掩码或自由答案，这些输出需要其他接口。对比学习的通用目标与采样问题可参见\BookVolumeRef{vol:paradigms}{第三卷“范式”}，这里的关键是图文两侧如何共同构成配对判定。

\subsection{ImageBind的图像配对与多模态表示联系}
\label{sec:vision-models-imagebind}

图文配对之外，视频帧和同时录到的声音、彩色图与对应深度、佩戴设备的画面与惯性信号，也提供了跨来源联系。若必须为所有模态两两采集配对数据，模态种类增加后组织成本会迅速上升。ImageBind用图像作为配对中心，把多个编码器接到已有视觉语义空间中\scite{vision-imagebind2023}

原实现分别编码图像、文字、音频、深度、热成像和惯性信号。图像与视频共用视觉编码器，视频输入由短时间范围内的帧构成。音频转换为梅尔频谱后划块，深度和热成像以单通道图像形式组织，惯性信号则把加速度计和陀螺仪的时间序列通过一维投影送入Transformer。各编码器末端有独立线性投影，输出归一化到同一维数。

以声音为例，一批训练数据包括画面$\bm{X}_i$与对应声音$a_i$。模型比较$\bm{v}_i$和$\bm{u}^{\mathrm a}_j$，沿图像到音频、音频到图像两个方向使用与式\eqref{eq:vision-clip-loss}同类的配对目标。深度、热成像和惯性数据各自与其配对画面计算目标，不要求一条记录同时拥有全部模态。原论文实现冻结已预训练的图像、文字编码器，学习音频、深度、热成像和惯性编码路径。

固定图像和文字空间的意义在于，新增模态需要向已经形成的坐标关系适配，而不是让所有编码器同时任意移动。训练仍需真实对应和足够数据覆盖。画面中既有车又有树而录音只有风声时，整段音画配对也没有显式指出声音来自哪个区域。

使用声音检索画面时，只需编码声音，再按它与图库向量的内积排序。使用声音选择文字类别时，可以比较声音向量与候选文字向量，即使新增音频编码器没有直接用音文对训练。论文在所测数据中观察到了这种未直接配对模态之间的联系；共同空间本身不是保证联系正确的定理。

一个几何事实有助于理解这种联系的可能性。若某段声音和某段文字都与同一个图像向量很接近，即$\lVert\bm{u}^{\mathrm a}-\bm v\rVert_2\leq\epsilon_{\mathrm a}$、$\lVert\bm{u}^{\mathrm t}-\bm v\rVert_2\leq\epsilon_{\mathrm t}$，三角不等式给出
\[
  \lVert\bm{u}^{\mathrm a}-\bm{u}^{\mathrm t}\rVert_2
  \leq\epsilon_{\mathrm a}+\epsilon_{\mathrm t}\text{。}
\]
但实际训练不能让每个语义一致的样本都围绕同一个准确图像锚点，也不能消除背景、视角和声源歧义，因此这只是有条件的几何解释。ImageBind的音频到图像生成展示还外接了图像生成系统，用声音表示替代相应条件；ImageBind本体仍输出表示。

\subsection{理解结果的形成与接口选择}
\label{sec:vision-models-understanding-interface}

街景查询“找相似画面”可以直接返回向量排序。“描述发生了什么”没有固定候选句表，需要根据视觉内容预测一串文字。“指出右侧车辆”则要求把表达对应到区域，仅有一句“右边有一辆车”尚未给出可检查的位置。三者都涉及理解，但需要的结果不同。

相似度适合在已定义候选中选择，语言生成适合组织可变长度的描述和答案，定位接口则要输出框、点或掩码。语言模型也可以按约定输出坐标，但仍需训练依据、坐标转换和定位评价。具体的语言指代与区域对应见\ref{sec:vision-understanding-grounding}节。

同一表示也可服务生成。图像语义向量可以作为生成器的条件，图文相似度的梯度可以调整采样方向，生成结束后的相似度还可以筛选候选。这些用途分别改变生成输入、生成过程和选择规则，不能混为“CLIP生成图像”。下一节从预测对象和训练模块出发，解释这些不同连接。

\section{多模态生成模型}
\label{sec:vision-models-generation}

生成描述、答案和语音响应，主要需要把已观察到的内容表达出来；生成图像、视频和环境声音，还要预测完整媒体中的大量细节。两类输出可以在一个系统中结合，所需目标和解码器却不同。本节分别追踪语言响应和媒体内容的生成路径。

\subsection{语言响应的生成}
\label{sec:vision-models-language-generation}

设观测经过编码得到视觉表示$\bm{F}$，请求为$q$，回答的词元为$y_1,\ldots,y_m$。条件语言生成通常按此前已经产生的前缀逐词展开：
\begin{equation}
  p_{\bm\theta}(y_1,\ldots,y_m\mid\bm F,q)
  =\prod_{j=1}^m p_{\bm\theta}(y_j\mid y_{<j},\bm F,q)\text{。}
  \label{eq:vision-vlm-language}
\end{equation}
训练时，把真实回答前缀提供给模型，以目标词元的负对数概率监督；使用时，前缀逐步由模型自己产生。式中的条件依赖只有通过实际连接才有可能成立。下面几种方法改变的重点，就是视觉信息怎样到达预测$y_j$所用的语言隐状态，以及这条路径上哪些权重被训练。

\subsubsection{BLIP-2的查询提取与视觉前缀}
\label{sec:vision-models-blip2}

图像编码器输出数百个特征，已有语言模型却没有学过怎样读取它们。BLIP-2在二者之间训练Q-Former，即用一组可学习查询从图像特征中抽取信息的Transformer。查询在这里是模型参数，不是用户的问题；其输出是一组随图片改变的连续表示\scite{vision-blip22023}

设图像特征为$\bm F\in\mathbb R^{n_{\mathrm v}\times d_{\mathrm v}}$，查询矩阵为$\bm Q_0\in\mathbb R^{r\times d_{\mathrm q}}$。查询通过自注意力交换信息，并以交叉注意力读取$\bm F$，形成$\bm Z\in\mathbb R^{r\times d_{\mathrm q}}$。原论文使用32个查询，无论图像编码器输出多少特征，交给下一模块的查询输出数量保持固定。压缩减少后续语言计算所读的视觉单元，也要求查询选择并保留与任务相关的信息。

第一阶段连接冻结图像编码器和Q-Former，使用图文对同时学习三种目标。对比目标让图像查询和文字分别编码，二者在自注意力中不能互看，以免把待匹配文本直接泄漏给图像表示；每段文字与多个查询输出比较，取最高相似度作为图文分数。配对判定目标则允许查询和文字双向交互，由各查询输出的二分类logit取平均，判断图文是否匹配，并使用困难负例。

第三种目标是图像条件文字生成。查询之间可以互看，但不能读取目标文字；每个文字词元可以读取全部查询和此前文字。由于文字不能直接读取冻结图像特征，生成目标迫使Q-Former经过查询路径传递有用信息。三种目标共用Q-Former参数，但注意力掩码不同，分别约束独立匹配、细致核验和信息表达。

第二阶段把$\bm Z$线性投影到语言模型的嵌入维数$d_{\mathrm l}$，得到$\bm P=\bm Z\bm W^\top\in\mathbb R^{r\times d_{\mathrm l}}$。这些向量放在文字输入前面，称为\term{视觉前缀}（visual prefix）：它们占据语言模型的输入位置，用连续向量提供图像条件，却不对应可直接打印的文字。冻结图像编码器和大型语言模型后，生成损失训练Q-Former与投影层。对于解码器式语言模型，目标是根据视觉前缀生成图像描述；对于编码器—解码器式语言模型，论文还使用前缀语言建模，把部分文字与视觉前缀交给编码器，预测剩余文字。

在街景问答中，模型先计算图像特征，再取得32个查询输出，接上“右侧停着什么车”及相应提示，由语言模型逐词生成答案。固定图片的视觉前缀可用于不同文字请求，但压缩时未保留的细节无法仅靠换一种问法恢复。这里的冻结说明针对两阶段预训练，不能推广到原论文所有下游微调设置，更不能作为所有VLM的固定训练顺序。

\subsubsection{Flamingo的视觉重采样与交错上下文}
\label{sec:vision-models-flamingo}

若输入包含“示例图像—示例回答—待答图像—问题”，视觉信息还需要与文字在序列中的位置对应。Flamingo先把每幅图像或视频压缩为固定数量的视觉单元，再让语言模型在中间层读取这些单元。原模型冻结已有视觉编码器与语言模型，训练新增连接部分\scite{vision-flamingo2022}

视觉编码器输出空间特征；视频则对采样帧分别编码，并加入时间嵌入。Perceiver Resampler使用可学习查询，通过注意力把数量可变的特征重采样为64个输出。这里的“重采样”是学习式信息汇聚，不是从原图均匀抽取64个像素。其注意力的键和值还包含查询隐表示本身，使这些潜在单元能够参与后续汇聚。

与BLIP-2把视觉表示放在语言输入前面不同，Flamingo在冻结的语言层之间插入门控交叉注意力与前馈模块。交叉注意力的查询来自当前语言隐状态，键和值来自重采样后的视觉表示。若用$C(\bm H,\bm Z)$概括一个新增交叉注意力分支，残差更新的核心形式为
\begin{equation}
  \bm H'=\bm H+\tanh(\alpha)\,C(\bm H,\bm Z)\text{，}
  \label{eq:vision-flamingo-gate}
\end{equation}
其中$\alpha$是从0开始学习的标量门控。初始时新增分支对输出的贡献为0，语言模型保留原来的计算；门控和连接权重随后从文字预测损失中学习。完整模块还包含相应的前馈分支，不应把式\eqref{eq:vision-flamingo-gate}当成整层结构。

交错输入还需要控制可见范围。原Flamingo的每个文字位置直接交叉注意到它之前最近的一幅图像或一个视频的视觉表示，而不是同时直接读取全部历史图片。历史图像仍可通过此前文字隐状态和语言自注意力间接影响结果。因此，示例中的红车和待答图中的红车虽然类别相同，也必须保持来源位置，不能把示例属性无条件移到目标上。

训练数据包含网页中的交错图文、图文对及视频文字对，目标仍是对可见视觉上下文条件下的文字进行预测。使用时，把几个已知示例和新请求放入上下文即可，示例本身不引起梯度更新。它们帮助说明当前要做什么任务，但示例数、输入处理和上下文压缩仍限制能保存的细节。图\ref{fig:vision-language-connectors}把前缀、层间读取与直接投影三条路径放在一起，便于区分视觉信息实际进入语言计算的位置。

\subsubsection{LLaVA的特征投影与视觉指令适配}
\label{sec:vision-models-llava}

若保留图像的网格特征，可以直接把每个视觉单元映射到语言嵌入空间，再让语言主干学习怎样使用这些输入。原始LLaVA采用这一连接方式，并用视觉指令数据训练回答行为。本节依据原论文v2中的线性投影设置，后来的多层感知机连接器不属于这里的模型\scite{vision-llava2023}

CLIP视觉编码器从图片取得网格特征$\bm F$，投影矩阵把各位置的特征转为与文字嵌入同宽的$\bm P$。这些视觉单元与问题词元一起进入自回归语言模型。投影并没有把每个图像块翻译成一个汉字或单词，而是调整连续特征的坐标，使语言层可以对它们计算注意力。

训练分为特征对齐和视觉指令适配。对齐阶段用图像描述构造“请简要描述图片”一类请求，冻结视觉编码器与语言模型，只更新投影层。语言模型的误差仍需沿计算图传回投影输入，冻结权重并不等于可以截断这条反向传播路径。指令适配阶段继续冻结视觉编码器，同时更新投影与语言模型，数据则包括对话、详细描述和复杂问答。

对多轮数据，输入保留图片、用户请求和前轮回答，损失监督助手应生成的回答词元。若请求是“比较两辆车的状态”，模型依赖同一视觉网格形成两个对象的描述，并逐词预测比较句。训练示例帮助模型学会“比较”这一请求形式，能否正确读取车的位置、朝向和遮挡仍由视觉证据及训练覆盖决定。

原模型可以生成自然语言，但没有因此得到任意精度的坐标、实例掩码或图像生成能力。若回答写了“右车正在行驶”，单张静态图片只提供姿态和场景线索，通常不足以确定实际速度；模型必须避免把语言中常见的情境补全当成观测。

\begin{figure*}[htbp]
  \centering
  % 每行仅保留用于区分接入位置的模块，不展开完整Transformer层。
\begin{tikzpicture}[x=1cm,y=1cm,font=\figurefont\small,
  block/.style={draw=FigureModel,fill=FigureModel!10,line width=0.8pt,
    minimum height=9mm,minimum width=21mm,align=center},
  input/.style={block,draw=FigureInput,fill=FigureInput!10},
  output/.style={block,draw=FigureOutput,fill=FigureOutput!10},
  arr/.style={-{Stealth[length=2mm]},BookInk,line width=0.8pt}]
  \node[anchor=west,font=\figurefont\small\bfseries] at (0,6.7) {BLIP-2：查询压缩后作为输入前缀};
  \node[input] (bf) at (1.1,5.7) {视觉特征};
  \node[block] (bq) at (4.1,5.7) {Q-Former};
  \node[block] (bp) at (7.2,5.7) {投影后的\\视觉前缀};
  \node[block] (bl) at (10.7,5.7) {语言模型};
  \node[output] (bo) at (14.1,5.7) {回答词元};
  \draw[arr] (bf) -- (bq);
  \draw[arr] (bq) -- (bp);
  \draw[arr] (bp) -- (bl);
  \draw[arr] (bl) -- (bo);
  \node[anchor=north] at (4.1,5.15) {可学习查询读取图像};
  \node (btext) at (10.7,4.7) {文字问题};
  \draw[arr] (btext) -- (bl);

  \node[anchor=west,font=\figurefont\small\bfseries] at (0,3.9) {Flamingo：在语言层之间读取视觉表示};
  \node[input] (ff) at (1.1,2.45) {视觉特征};
  \node[block] (fr) at (4.1,2.45) {Perceiver\\Resampler};
  \node[block] (fc) at (7.2,2.45) {门控交叉\\注意力};
  \node[block] (fl) at (10.7,2.45) {后续语言层};
  \node[output] (fo) at (14.1,2.45) {回答词元};
  \draw[arr] (ff) -- (fr);
  \draw[arr] (fr) -- node[above] {$K,V$} (fc);
  \draw[arr] (fc) -- (fl);
  \draw[arr] (fl) -- (fo);
  \node (fh) at (7.2,3.35) {语言隐状态$Q$};
  \draw[arr] (fh) -- (fc);
  \node[anchor=north] at (4.1,1.9) {每个图像／视频64个单元};

  \node[anchor=west,font=\figurefont\small\bfseries] at (0,0.7) {原始LLaVA：网格特征投影后进入共享序列};
  \node[input] (lf) at (1.1,-0.25) {视觉网格};
  \node[block] (lp) at (4.1,-0.25) {线性投影};
  \node[block] (ls) at (7.2,-0.25) {视觉单元\\与问题序列};
  \node[block] (ll) at (10.7,-0.25) {语言模型};
  \node[output] (lo) at (14.1,-0.25) {回答词元};
  \draw[arr] (lf) -- (lp);
  \draw[arr] (lp) -- (ls);
  \draw[arr] (ls) -- (ll);
  \draw[arr] (ll) -- (lo);
\end{tikzpicture}

  \caption{三种视觉接入方式的结构示意，按各原论文的预训练路径绘制。BLIP-2由少量查询形成视觉前缀，Flamingo在语言层中新增读取视觉表示的交叉注意力，原始LLaVA把视觉网格线性投影后交给语言主干。图中省略各模块内部层，框旁文字指出视觉信息进入的位置；冻结范围由正文所述阶段决定。}
  \label{fig:vision-language-connectors}
\end{figure*}
\FloatBarrier

\subsubsection{Qwen3-VL的多层视觉接入与视频时间组织}
\label{sec:vision-models-qwen3vl}

多页文档和视频包含分散证据，只读取视觉编码器的最末层、再把所有帧当成无时间含义的图片，可能丢失细节与先后关系。Qwen3-VL技术报告v1在动态分辨率编码、视觉合并、跨层接入和时间表示上组织这些信息。报告区分思考版与非思考版，以下机制不意味着两类后训练版本具有相同响应行为。

视觉编码器接受动态分辨率输入，输出长度随图像大小变化的视觉特征。视觉语言合并模块使用两层MLP，把相邻$2\times2$特征压缩为一个语言宽度的视觉单元\scite{vision-qwen3vl2025}
空间合并降低上下文长度，同时也把四个位置的信息汇集到一个位置表示中；实际细节保留需结合输入分辨率和训练检查。

DeepStack进一步利用视觉编码器三个不同深度的中间特征。每组特征经独立合并模块映射后，加到语言主干前三层相应的视觉位置隐状态上。视觉信息由多个深度进入语言计算，而不是全部仅依赖初始输入。残差接入不增加额外序列位置，但仍增加了特征投影和存储等计算。

位置表示有两项不同安排。交错MRoPE把时间、高度、宽度对应的旋转位置维度分散到不同频率范围，避免把某个坐标维度集中在单一频段；这是注意力内部如何表示相对位置的问题。视频时间戳则以文字形式放在每个时间块之前，例如“3.0秒”，提供明确的时间读数。两者共同使用，但文字时间戳不是从MRoPE隐向量中自动打印出来的数值。

在一个教学视频中，采样块对应1.0、3.0和8.0秒。保留这些时间戳，模型才有条件判断后两个观测相隔5秒；只给顺序编号会把这种间隔信息删掉。不过，如果敲门动作发生在4.2秒且没有进入任何采样块，时间戳再准确也无法补回动作。多图输入同样要保留图号，避免把第二幅图的车窗状态归到第一幅图。

报告的预训练先只更新合并模块，冻结视觉编码器和语言主干，随后转为全部模块训练，并逐步扩展上下文窗口。后训练包含指令适配、蒸馏与强化学习等安排，思考版使用相应的推理形式数据，非思考版使用标准回答形式。使用时应选定具体检查点、格式和观察预算。报告支持的256K词元上下文是容量约定，不是对任意长文档小字和短暂视频事件的保存保证。

\subsubsection{Qwen2.5-Omni的音画输入与文字语音响应}
\label{sec:vision-models-qwen-omni}

视频中的敲门动作和敲击声可能互相补充，也可能错位或来自画外。Qwen2.5-Omni技术报告v1通过Thinker--Talker结构处理这类输入并输出文字与语音。Thinker读取多模态上下文、生成文字及隐表示，Talker利用这些信息生成语音编码，再恢复成声音\scite{vision-qwenomni2025}

输入音频被转换为梅尔频谱，经音频编码器形成时间特征；图像和视频由视觉编码器处理。时间对齐多模态旋转位置编码TMRoPE给音视频提供共同时间尺度，报告中一个时间位置单位对应40毫秒。音频单元沿时间排列，视频单元同时保留帧时间和空间位置。输入还按实际时间划分为2秒块，每块内先放视觉表示、后放音频表示，再按块交错。

因此，序列中相邻的视觉和声音单元不一定只相差一个物理采样点，实际时间由位置记录共同表达。这个安排也不证明两段内容来自同一个声源。若门外另有人敲击金属，画面中的手部动作与声音可能在时间上接近，却缺乏确定来源。回答应分别依据看见的动作、听见的声音和额外文字线索。

Thinker按式\eqref{eq:vision-vlm-language}的同类方式预测回答文字。Talker接收Thinker的高层隐表示以及已采样文字词元的嵌入，以双轨自回归方式组织文字和语音编码。隐表示提供语义、语气等上下文，离散文字则帮助确定具体词语，避免语义接近但读音不同的内容仅由连续表示区分。

语音编码还不是可以播放的波形。报告的流式解码路径用流匹配模型将编码转为梅尔频谱，再由声码器合成波形；\ref{sec:vision-models-flow-matching}节将解释所需的速度预测目标。解码采用有限块上下文，包含回看和一定前看，因此流式表示分块产生结果，并不意味着没有等待。

训练职责也有区别。多模态预训练起步时固定语言模型，先适配连接再训练相关编码器，随后展开全参数与长上下文训练；Thinker再用多模态指令数据适配回答。Talker单独经历语音续接、偏好优化和说话人适配等训练安排，以学习从语义上下文到语音的生成。理解监督不能替代语音发音监督，音画时间对齐也不能替代输出声学解码。

在敲门例子中，输入路径保留动作和声音，Thinker形成“画面中的人抬手敲门，并伴有敲击声”这样的候选回答，Talker把确定的回答组织为语音。该版本支持文字与语音输出，不因此具有图像或视频生成模块。首段声音的到达时间还包含输入编码、文字生成、语音编码预测及波形解码等开销。

\subsection{媒体内容的生成}
\label{sec:vision-models-media-generation}

生成一幅雨天街景时，仅有“雨天”这个语义判断还不够，模型需要决定每处颜色、纹理与空间关系。直接在像素中生成可以保留细粒度自由度，但计算量大；在压缩表示中生成可以减少计算，却必须有相应解码器。条件信息怎样进入生成过程，又形成另一项独立选择。

\subsubsection{LDM的空间潜表示与文字条件接入}
\label{sec:vision-models-ldm}

潜空间扩散模型（latent diffusion model, LDM）先学习图像压缩与重建，再在压缩空间中学习生成。这里的潜表示保留二维网格结构，与CLIP把整图汇总成一个语义向量的空间不同。原论文的文字条件LDM通过交叉注意力把文字表示送入去噪网络\scite{pgm-generative-rombach2022}

图像自编码器由编码器$E$与解码器$D$组成。编码得到$\bm Z=E(\bm X)\in\mathbb R^{h'\times w'\times d_z}$，解码得到$D(\bm Z)$。自编码阶段使用重建相关的感知与对抗目标，并配合潜空间正则化，学习在压缩后保留可重建的信息。该阶段完成后固定自编码器，在其潜空间中训练扩散网络，不把图像重建器和每个生成任务从头共同重训。

为说明计算对象，把潜数组按固定次序展开为向量$\bm z$。取噪声步$k$，给它加入标准高斯噪声$\bm\epsilon$：
\[
  \bm z^{(k)}
  =\sqrt{\bar\alpha_k}\bm z+
    \sqrt{1-\bar\alpha_k}\bm\epsilon\text{，}
\]
其中$\bar\alpha_k$由噪声日程指定，$(k)$表示加噪或去噪过程的步，不是视频帧。文字请求$c$经条件编码器得到$\bm C=E_{\bm\phi}(c)$。去噪网络以当前潜表示、噪声步和文字为条件预测所加噪声，目标为
\begin{equation}
  L_{\mathrm{LDM}}
  =\mathbb E_{\bm X,c,k,\bm\epsilon}
    \left[
      \left\lVert\bm\epsilon-
      \epsilon_{\bm\theta}
      \bigl(\bm z^{(k)},k,E_{\bm\phi}(c)\bigr)
      \right\rVert_2^2
    \right]\text{。}
  \label{eq:vision-ldm-objective}
\end{equation}
原论文的条件设置共同优化$\bm\theta$与$\bm\phi$；后来使用冻结文字编码器的实现，需要另行说明。

交叉注意力中，查询来自当前去噪网络的空间特征，键和值来自文字序列。假设当前有$n_z$个空间单元、文字有$m$个词元，注意力权重是$n_z\times m$的矩阵，每个空间位置可以按当前状态读取不同文字信息。它没有直接为“红色”指定一个真实对象框，词语与空间的正确对应仍需通过数据和目标学习。

使用时先编码“雨天街道，右侧停着红色汽车”，再从潜空间噪声出发多步去噪，最终用$D$恢复图片。若高宽都压缩为原来的八分之一，潜网格位置数变为$1/64$；这并不意味着全部计算或内存恰好减少64倍，因为通道数、网络宽度和注意力实现仍然不同。压缩还可能损失小字等细节，生成后的事实与控制要求由\ref{chap:vision-image-generation}章的任务协议检查。

\subsubsection{unCLIP的图像表示先验与图像解码}
\label{sec:vision-models-unclip}

LDM直接用文字条件预测空间潜表示，unCLIP则在文字与图像之间加入一个整图语义表示。冻结CLIP后，先验模型根据文字预测CLIP图像表示，另一个生成解码器再依据该表示形成图片。这是DALL-E 2报告中的unCLIP路线，生成器并非CLIP本身\scite{vision-unclip2022}

设文字为$c$，CLIP图像表示为$\bm z_{\mathrm i}$，图像为$\bm X$，模型可写成
\begin{equation}
  p(\bm X\mid c)
  =\int p_{\mathrm{dec}}(\bm X\mid\bm z_{\mathrm i},c)\,
          p_{\mathrm{prior}}(\bm z_{\mathrm i}\mid c)\,
          \mathrm d\bm z_{\mathrm i}\text{。}
  \label{eq:vision-unclip-factorization}
\end{equation}
这里表达模型的分层生成分布，不宣称训练后已等于真实数据分布。先验处理“这段文字可能对应哪些图像语义表示”，解码器处理“这个语义表示可能对应哪些具体像素”。后者还可以读取文字，不能把完整条件一律删成只有图像向量。

训练先验时，图文对中的图像经冻结CLIP形成目标向量。论文比较了两条路线：自回归先验先对图像向量作主成分降维与分量量化，再逐项预测离散编码；扩散先验直接对连续图像向量加噪，并在文字条件下回归干净向量。扩散先验采用干净表示预测目标，不能直接照搬LDM的噪声预测符号而不作转换。

训练生成解码器时，以真实图像的CLIP表示和相应文字为条件，对带噪图像学习去噪。图像表示被投影到时间步嵌入及额外上下文单元，连接到由GLIDE结构改造的图像生成网络；随后使用图像上采样模块恢复更高分辨率。CLIP在先验与解码器训练期间保持冻结，两个生成模块学习不同的映射。

使用文字时，先采样语义表示，再采样图像。使用参考图片生成变体时，可以直接编码该图片，跳过文字到图像表示的先验预测。由于整图语义表示没有保留全部像素，不同采样能产生语义接近而纹理、布局细节不同的结果。若要求准确重建某张图片，还需保留解码反演等额外信息；“语义相似”不能替代对原图每处细节的保持。

\subsubsection{GLIDE的文字条件与语义引导}
\label{sec:vision-models-glide}

条件输入规定去噪网络能够读取什么，引导规定使用时怎样组合或修正预测，生成后重排则只决定从已有候选中选谁。GLIDE提供了比较这三个环节的具体入口。图像生成章的\ref{sec:vision-generation-glide-guidance}节已展开采样公式，本节把它们放回模块关系中。

GLIDE先用文字Transformer编码请求，将文字信息接入图像去噪网络。训练根据图文配对构造带噪图像，让去噪预测依赖文字。为了实现无分类器引导，训练时还以一定概率去掉文字条件，使同一个网络能给出条件预测和无条件预测。使用时对同一噪声状态分别计算两者，再按选定强度组合。

噪声感知CLIP引导需要另一套图文比较模型。它必须读取当前噪声水平的图像，把文字相似度对图像的梯度用于修正生成方向。原GLIDE实验专门训练了适应噪声输入的CLIP，不能把只看清晰图片的原CLIP不加说明地放在同一位置。固定请求与初始噪声时，改变引导强度会改变每一步的预测或更新，最终像素也可能随之改变。

生成后重排则先完成$K$个候选，之后才计算各自的匹配分数并选择。重排不修改任何候选，也不能从全部错误的候选中创造正确答案。三种环节可以组合，但应把额外模型调用、梯度和候选数量计入预算。图文相似度仍需与对象数量、属性归属和空间关系的独立核验配合。

\subsubsection{流匹配的速度目标与连续生成}
\label{sec:vision-models-flow-matching}

媒体生成不一定把网络输出解释为“应去掉多少噪声”，也可以直接学习“当前表示应向哪个方向变化”。\term{流匹配}（flow matching）就是通过回归预先构造路径上的速度来训练连续生成模型。它把训练时可计算的局部目标，与使用时求解整条生成轨迹联系起来\scite{vision-flowmatching2023}

用一个明确的教学构造说明。设数据潜向量为$\bm z_1\in\mathbb R^d$，噪声为$\bm z_0\sim\mathcal N(\bm0,\bm I)$，路径参数$s\in[0,1]$。在两者之间定义直线路径
\begin{equation}
  \bm z(s)=(1-s)\bm z_0+s\bm z_1,
  \qquad
  \bm u=\frac{\mathrm d\bm z(s)}{\mathrm ds}
       =\bm z_1-\bm z_0\text{。}
  \label{eq:vision-flow-path}
\end{equation}
$s=0$对应噪声，$s=1$对应数据。这里采用无残余噪声的理想化端点以便手算；原流匹配论文也讨论终点保留微小方差等路径，不能把所有实现的日程等同于这一个式子。

训练时取一对数据与条件$(\bm z_1,c)$，抽样噪声$\bm z_0$和时刻$s$，就能直接算出中间状态$\bm z(s)$与目标速度$\bm u$。速度网络$\bm v_{\bm\theta}$的平方误差目标为
\begin{equation}
  L_{\mathrm{FM}}
  =\mathbb E_{\bm z_1,c,\bm z_0,s}
     \left[
       \left\lVert
         \bm v_{\bm\theta}(\bm z(s),s,c)
           -(\bm z_1-\bm z_0)
       \right\rVert_2^2
     \right]\text{。}
  \label{eq:vision-flow-loss}
\end{equation}
计算一次训练目标不必先求解从0到$s$的微分方程，因为中间状态有显式构造。更新的是速度网络参数，样本端点和抽样噪声给出监督。

使用时却不知道目标$\bm z_1$，所以不能直接取$\bm z_1-\bm z_0$走向答案。模型需要从一个新噪声出发，根据当前状态预测速度，求解
\begin{equation}
  \frac{\mathrm d\widehat{\bm z}(s)}{\mathrm ds}
    =\bm v_{\bm\theta}(\widehat{\bm z}(s),s,c),
  \qquad
  \widehat{\bm z}(0)\sim\mathcal N(\bm0,\bm I)\text{，}
  \label{eq:vision-flow-ode}
\end{equation}
再把$\widehat{\bm z}(1)$送入媒体解码器。简单的欧拉法在$s_k$处用
$\widehat{\bm z}^{(k+1)}
=\widehat{\bm z}^{(k)}+(s_{k+1}-s_k)
\bm v_{\bm\theta}(\widehat{\bm z}^{(k)},s_k,c)$
更新。有限步长有数值误差，网络也有预测误差，两者不能混为同一项。

为什么从许多不同样本路径学习，仍能得到一个只依赖当前状态的速度？固定输入$(\bm z,s,c)$时，平方损失的最优预测是目标速度的条件均值。令$\bar{\bm u}=\mathbb E[\bm u\mid\bm z,s,c]$，则
\[
  \mathbb E[\lVert\bm v-\bm u\rVert_2^2\mid\bm z,s,c]
  =\lVert\bm v-\bar{\bm u}\rVert_2^2+
    \mathbb E[\lVert\bm u-\bar{\bm u}\rVert_2^2\mid\bm z,s,c]\text{。}
\]
展开平方后，交叉项因$\mathbb E[\bm u-\bar{\bm u}\mid\bm z,s,c]=\bm0$消失。右侧第二项不依赖$\bm v$，因此网络学习的是穿过当前状态的样本路径速度的条件平均。在适当正则条件下，这一边缘速度场可以推动对应的概率分布路径；它不要求网络辨认训练时某条唯一的端点连线。

取二维教学样本$\bm z_0=(0,1)^\top$、$\bm z_1=(2,-1)^\top$，目标速度为$(2,-2)^\top$，$s=0.25$时状态为$(0.5,0.5)^\top$。假如使用时在该位置预测速度$(1.2,-2.8)^\top$，步长0.25给出下一状态$(0.8,-0.2)^\top$，而该训练直线在$s=0.5$处为$(1,0)^\top$。这个差异说明一次速度预测误差怎样传到位置，不能用单条训练直线宣称网络推理必然沿直线运行。图\ref{fig:vision-flow-step}给出几何对应。

\begin{figure}[htbp]
  \centering
  % 直线路径为教学构造；粗箭头使用特意偏离目标的速度预测。
\begin{tikzpicture}[x=3cm,y=1.65cm,font=\figurefont\small,
  arr/.style={-{Stealth[length=2mm]},BookInk,line width=0.8pt}]
  \draw[arr] (-0.15,-1.4) -- (2.55,-1.4) node[right] {$z_1$};
  \draw[arr] (-0.15,-1.4) -- (-0.15,1.5) node[above] {$z_2$};
  \foreach \x in {0,1,2}{
    \draw[BookMuted,line width=0.5pt] (\x,-1.4) -- (\x,-1.46);
    \node[below] at (\x,-1.46) {$\x$};
  }
  \foreach \y in {-1,0,1}{
    \draw[BookMuted,line width=0.5pt] (-0.15,\y) -- (-0.19,\y);
    \node[left] at (-0.19,\y) {$\y$};
  }
  \draw[BookBlue,line width=0.9pt] (0,1) -- (2,-1);
  \fill[BookBlue] (0,1) circle[radius=1.1mm];
  \fill[BookTeal] (2,-1) circle[radius=1.1mm];
  \node[anchor=south west] at (0,1.13) {$s=0,\ (0,1)$};
  \node[anchor=west] at (2.02,-0.96) {$s=1$};
  \fill[BookInk] (0.5,0.5) circle[radius=1mm];
  \node[anchor=south west] at (0.52,0.62) {$s=0.25,\ (0.5,0.5)$};
  \draw[BookTeal,line width=1pt,fill=white] (1,0) circle[radius=1.1mm];
  \node[anchor=west] (ideal) at (1.35,0.1) {路径上$(1,0)$};
  \draw[BookTeal,line width=0.6pt] (1.04,0) -- (1.3,0.1);
  \draw[-{Stealth[length=2.2mm]},BookAmber,line width=1.3pt]
    (0.5,0.5) -- (0.8,-0.2);
  \fill[BookAmber] (0.7825,-0.23) rectangle (0.8175,-0.17);
  \node[anchor=east] at (0.76,-0.35) {预测$(0.8,-0.2)$};
  \draw[BookMuted,dashed,line width=0.8pt] (0.82,-0.17) -- (0.97,-0.03);
  \node[anchor=west] at (0,-2.05) {粗箭头：$0.25(1.2,-2.8)^\top$};
\end{tikzpicture}

  \caption{流匹配教学样本的路径与一次欧拉更新。实线连接已知训练端点，粗箭头为预测速度乘以步长后的位移，虚线标示与训练路径上$s=0.5$位置的差别。横纵轴是两个潜变量坐标；路径参数$s$不是视频中的事件时间。}
  \label{fig:vision-flow-step}
\end{figure}

即便每条条件路径都很简单，整体数据分布对应的轨迹也未必是直线，更不能仅据此声称得到了全局最优传输。规范化流描述通过可逆变换建立分布的模型，流匹配描述训练速度场的一类目标，两者相关但不是同义词。需要连续密度与变换公式时，可回查\BookSectionRef{sec:pgm-generative-flows}{第二卷“生成模型”章的规范化流一节}。

\subsubsection{LTX-2的音视频双分支与联合解码}
\label{sec:vision-models-ltx2}

若只生成敲门画面，再由另一个模型补声音，后一步只能适应已固定的动作。联合生成可以让画面和声音在生成过程中互相影响。LTX-2技术报告v1采用分别表示、相互连接的音视频分支，学习文字条件下的联合媒体输出。

视频经因果VAE编码为空间与时间潜网格，音频经频谱转换和独立音频VAE编码为时间潜序列。两个分支的长度、通道宽度和压缩方式可以不同，因此不需要把音频潜向量误解释成视频像素。文字经编码与连接模块形成条件，分别进入两种媒体的生成分支。

每个分支具有自己的注意力和前馈计算，又通过双向交叉注意力交换信息。视频查询读取音频的键和值，音频查询反过来读取视频的键和值。视频分支内部同时使用时空位置，跨音视频连接则着重使用时间位置，因为音频时间单元没有图像平面上的横纵坐标。时间编码使对应关系可被表达，但正确的敲击同步仍要从配对数据和生成结果检验。

训练把配对音画编码为各自目标潜表示，构造噪声状态，用流匹配监督两个分支的速度\scite{vision-ltx22026}
沿用前节的教学记号，可以把联合状态写成$(\bm z^{\mathrm v}(s),\bm z^{\mathrm a}(s))$，两边的速度预测都依赖当前音画状态及文字条件。各项损失的尺度需要按维度和具体训练设置处理，不能不计音视频长度差异而把元素误差总和直接比较。

使用时，从两个潜空间的噪声开始共同更新。视频潜表示由视频VAE解码为帧，音频潜表示先恢复梅尔频谱，再由声码器恢复波形。若文字要求“人物敲门两次，随后停下”，验收至少要对应两次接触、两次声音以及停顿，不能只听到类似敲门声就认定通过。联合生成不同于为固定视频配声，接口本身也不提供视觉问答；LTX-2中的文字编码器服务生成条件，不等于系统具有自由问答输出头。

\section{统一理解与生成模型}
\label{sec:vision-models-unified}

前面的方法可以形成可比较表示、语言或媒体。如果用户希望“解释这张图，再按要求画一个变体”，可以在一个主干中共享计算，也可以把理解模型和生成器连接起来。判断“统一”的含义，需要看哪些参数、表示和训练目标实际共用，不能只看界面是否只有一个输入框。

\subsection{Show-o2的共享视觉表示与双输出头}
\label{sec:vision-models-showo2}

理解图像需要语义特征，生成图像还需要纹理和局部结构。Show-o2在同一视觉潜空间中提取这两类信息，送入共享语言主干，再用不同输出头分别处理文字和视觉生成。本节依据论文v3的结构及训练说明\scite{vision-showo22025}

图像和视频由3D因果VAE编码。因果性约束编码中未来信息的使用，图像可以作为相应的单帧情形处理。视觉潜表示经过两条路径：语义层从潜网格中提取高层信息，投影层保留便于后续处理的低层特征。语义层在主要两阶段训练之前，利用预训练SigLIP的图像特征进行蒸馏，使其能够从干净和带噪潜表示中提取语义。

两个路径按对应的空间及时间位置对齐，沿特征维拼接，再经归一化与MLP融合为共享视觉表示。这里的共享表示保留多个位置，并非把整幅图压成一个CLIP向量。干净参考图以路径终点$s=1$的状态进入，生成中的视觉块则带有当前路径时刻。相同的视觉接入模块因此需要处理不同噪声水平。

文字嵌入和视觉块按交错顺序进入主干。序列块之间遵守因果顺序，使后续输出可以依赖先前条件；同一视觉块内部允许完整注意力，使一个待生成图像或视频块中的位置共同参与当前速度预测。这里的“因果”描述输出块的条件顺序，不等于每个视频帧只能读取此前帧。

语言头把隐状态转换为下一词元分布，流头则预测视觉潜表示的速度。后者包含带生成时刻调制的Transformer层，其输出经过数值更新和VAE解码才成为媒体。两种训练目标可以概括为
\begin{equation}
  L=\lambda_{\mathrm{text}}L_{\mathrm{NTP}}+L_{\mathrm{FM}}\text{，}
  \label{eq:vision-showo-loss}
\end{equation}
其中$L_{\mathrm{NTP}}$是文字的下一词元预测损失，$L_{\mathrm{FM}}$是视觉速度误差，$\lambda_{\mathrm{text}}$控制文字项权重。两项可以通过共享主干联系起来，但输出维数和误差对象不同。

原论文的第一阶段只训练投影、时空融合和流头，保留已有语言知识，并学习视觉生成接口；第二阶段用高质量理解与生成数据适配除VAE外的整体模型。图\ref{fig:vision-showo-heads}把共享范围与两个输出头分开表示。VAE不随这两个阶段更新，统一训练也不意味着每个部件都在每个阶段改变。

\begin{figure*}[htbp]
  \centering
  \begin{tikzpicture}[x=1cm,y=1cm,font=\figurefont\small,
  block/.style={draw=FigureModel,fill=FigureModel!10,line width=0.8pt,
    minimum height=9mm,minimum width=20mm,align=center},
  input/.style={block,draw=FigureInput,fill=FigureInput!10},
  output/.style={block,draw=FigureOutput,fill=FigureOutput!10},
  arr/.style={-{Stealth[length=2mm]},BookInk,line width=0.8pt,
    rounded corners=1.2mm}]
  \node[input] (z) at (1.1,1.8) {干净／带噪\\视觉潜表示};
  \node[block] (sem) at (4.0,2.7) {语义层};
  \node[block] (low) at (4.0,0.9) {低层投影};
  \node[block] (fuse) at (6.9,1.8) {按位置拼接\\与融合};
  \node[block] (trunk) at (9.8,1.8) {共享主干};
  \node[block] (text) at (12.7,2.7) {语言头};
  \node[block] (flow) at (12.7,0.9) {流头};
  \node[output] (words) at (15.4,2.7) {回答词元};
  \node[output] (media) at (15.4,0.9) {更新潜状态\\VAE解码};
  \draw[BookInk,line width=0.8pt] (z.east) -- (2.55,1.8);
  \draw[arr] (2.55,1.8) |- (sem.west);
  \draw[arr] (2.55,1.8) |- (low.west);
  \draw[arr] (sem.east) -- (5.45,2.7) |- (fuse.west);
  \draw[arr] (low.east) -- (5.45,0.9) |- (fuse.west);
  \draw[arr] (fuse) -- (trunk);
  \draw[BookInk,line width=0.8pt] (trunk.east) -- (11.25,1.8);
  \draw[arr] (11.25,1.8) |- (text.west);
  \draw[arr] (11.25,1.8) |- (flow.west);
  \draw[arr] (text) -- (words);
  \draw[arr] (flow) -- (media);
  \node (prompt) at (9.8,0.05) {文字上下文};
  \draw[arr] (prompt) -- (trunk);
  \node[anchor=west] at (0,-0.75) {视觉块之间保持因果顺序，块内位置允许完整注意力。};
\end{tikzpicture}

  \caption{Show-o2的共享视觉接入与双输出头示意，依据论文v3重绘。视觉潜表示的语义路径和低层路径按位置融合，再与文字一起进入共享主干。语言头给出词元，流头给出视觉潜变量的速度；生成分支需迭代更新潜状态并重新经过模型，图中省略迭代回边以突出共享和专用模块的边界。}
  \label{fig:vision-showo-heads}
\end{figure*}
\FloatBarrier

对于“解释参考图，再生成雨天变体”，参考图先经过干净视觉路径，语言头生成解释；随后的视觉块依据此前上下文，通过流头逐步形成媒体。解释正确并不自动保证生成正确，生成中仍可能丢失刚才已识别的对象关系。论文包含文字、图像和视频路线，但视频与交错数据的训练安排随变体而异，报告中的较大7B变体没有沿用同样的视频生成训练数据。选择检查点时必须确认实际训练范围，不能把模型族的全部展示合并成每个版本的能力，更不能额外推断音频生成。

\subsection{理解模型与专用生成器的接口组合}
\label{sec:vision-models-composition}

另一条路线是由视觉语言模型读取图片和请求，形成明确的内容说明，再调用专用生成器。以“把玩具狗旁的花瓶改成蓝色，其余尽量保持”为教学请求，理解模型可以输出目标短语“花瓶”、保留对象“玩具狗”和修改属性“蓝色”。若生成器支持局部编辑，还需要把目标短语转换成原图坐标中的区域或掩码，并把原图作为独立条件传递。

只传一句“蓝色花瓶旁有玩具狗”时，生成器没有收到原图的精确纹理和位置，因而只能依据语义重建场景。传原图与掩码可以保留更多空间条件，但若理解模型把目标区域选成玩具狗，生成器仍可能准确执行了错误区域上的请求。组合系统应保存中间说明、区域来源、输入图片和最终结果，使错误能定位到读取、传递还是生成环节。

连续特征也可以通过适配模块传给生成器，减少必须把全部细节说成文字的限制。然而，不同编码器的特征坐标和尺度未必兼容，适配模块通常需要自己的训练依据。分别训练后直接连接、固定两端学习适配器、共同更新多个模块，是不同的训练安排，不能仅因推理时连续调用就称为共同学习。

与Show-o2相比，这类组合通常不共享理解与生成主干，可以独立替换生成器或读取模型，却增加了中间格式转换和多次调用。共享主干减少某些重复计算，但也需要在同一组参数中协调任务。两种路线的选择取决于能传递哪些证据、哪一步容易出错，以及最终成果是否满足请求。

\section{模型能力与使用约束}
\label{sec:vision-models-constraints}

一个模型支持的输入和输出，是开始选型的依据。实际使用还需确认信息经过编码、压缩和生成后，是否保留了当前任务所需的内容，以及达到这一结果用了多少资源。下面从信息、训练依据和成本三个方面回到同一街景与音画请求。

\subsection{信息保留与输出能力}
\label{sec:vision-models-information}

整图缩小会让车牌字符变成几个像素，局部裁剪会删除周围关系，视频抽帧会漏掉短动作，固定查询数会把许多视觉位置汇总成少量表示。投影和共享主干可以重新组织已经保留的信息，却不能保证恢复这些阶段丢弃的证据。

可以固定问题，分别提供整图、局部高分辨率图，以及带原图位置的二者组合。如果只有局部图能正确读取车牌，就应进一步判断错误来自原始分辨率、输入缩放还是后续压缩，而不是立即归为语言推理失败。组合输入还要确认模型知道局部图来自哪里，否则可能把两个视图当成两辆不同的车。

输出也需逐项验收。匹配分数正确只说明当前候选中的选择，语言流畅只说明形式自然，视频画面逼真只说明感知质量的一部分。数量、位置、时间与来源需要对应的检查。精确区域和答案证据见\ref{chap:vision-multimodal-understanding}章，媒体连贯性和控制见\ref{chap:vision-multimodal-generation}章。

\subsection{模态贡献与训练依据}
\label{sec:vision-models-modality-contribution}

模型接受声音，并不说明一次回答确实使用了声音。可以在保持问题不变时移除音频，或换成时间错开的声音，检查结果与正确答案怎样变化。若问题只问车的颜色，移除声音后答案不变是合理的；若问题问敲击发生了几次，错误声音改变了答案，则需要判断模型是否把它与画面进行了适当核验。

同样，替换街景图却保留文字问题，可以检查模型是否依赖常见语言模式猜测。替换条件必须使正确答案或证据关系发生可解释的变化；只观察回答变了，无法证明模型使用了正确模态。多模态评价还需要覆盖缺失模态、冲突模态和训练中未直接出现的组合。

训练依据决定了可以期待哪一层联系。整图配对监督适合建立全局对应，框或时间区间标注提供更细的位置联系，指令示例说明请求和回答怎样组织。它们可以联合使用，但不能互相冒充。某个新组合的成功应在同类样本与清楚的任务协议中复核，不能由一次演示推断任意模态都已对齐。

\subsection{计算资源与使用成本}
\label{sec:vision-models-cost}

多模态成本可以按输入编码、联合上下文计算、输出生成和媒体解码记录。固定图像大小还不够，裁剪数、视频采样、音频长度、历史轮数和输出时长都可能改变实际输入与输出规模。一个只生成短答案的调用，不能与生成数秒音视频的调用仅比较总延迟。

在标准稠密注意力中，$n$个序列位置的成对分数规模为$n^2$，隐藏宽度为$d$时，相关乘加量含$O(n^2d)$项。若有256个文字单元和1024个视觉单元，总长为1280；视觉单元增加到4096时，总长4352，成对位置数约变为原来的11.56倍。这个教学比例只描述该计算项，不能直接当作实际延迟倍率，局部注意力、缓存与分块实现都会改变资源安排。

冻结模块减少需要保存的参数梯度和优化器状态，却不省去推理中的前向计算。若可训练连接器位于冻结语言模型之前，训练还可能需要经过冻结层计算输入梯度。媒体生成则通常反复调用速度或去噪网络，成本取决于步数和每步条件计算；候选重排还增加整条生成路径的重复次数。

共享主干和模型组合也需按实际调用比较。前者可能复用上下文，后者可以独立缓存编码或选择不同规模的专用模块。通用部署与资源管理可参见\BookVolumeRef{vol:systems}{第五卷“系统”}。对当前多模态问题，更有针对性的研究方向是怎样在固定预算中保留任务必需的局部证据、建立可靠时间对应，并协调不同输出目标；扩大名义容量只是其中一种手段。

\section{本章小结}
\label{sec:vision-models-summary}

街景照片、请求和声音进入同一系统，需要经过有来源与位置约定的编码，再由匹配空间、视觉读取或共享主干建立联系。CLIP与ImageBind说明可比较表示怎样形成；视觉前缀、层间交叉注意力和多层接入说明语言生成怎样利用视觉；潜空间生成与速度预测则为媒体输出提供了不同的计算路径。

这些路径可以在同一模型中共享，也可以通过专用模块组合。语言回答依靠下一词元预测表达判断，图像和音视频输出还需相应潜表示、生成更新与解码器。输入支持范围、训练过的对应关系和实际输出能力应分别确认。选好接口后，仍需回到任务检查答案与证据、内容与条件是否一致，这也是后续理解与生成两章的重点。
